\chapter{Architecture for reversible review}
\label{ch:architecture}

Open a finished review packet and four things should be easy to find: the
version the reader saw, the changes the author made, the technical objects that
survived, and the reader's decision. The packet exists only because an earlier
brief, blueprint, and pre-human release record cleared the path. Parsers,
linters, models, and build tools support those steps; none of them owns the
document.

That practical test drives the engineering. The system records explicit
correspondence and abstains when it cannot establish one. The author sees
context and alternatives. The reader sees a clean document and a short task.

Figure~\ref{fig:record-model} shows the minimum set of relationships. A larger
schema would add fields, but it would not change the ownership of the arrows.

\hvFigureRecordModel

\section{The objects that must survive a revision}
\label{sec:architecture-objects}

The first release needs eleven records. Their names are technical
because the software stores them, not because the finished document should
expose them as labels.

\begin{description}[style=nextline,leftmargin=37mm,labelindent=0mm]
\item[Authoring brief] The intended reader, decision, genre, known vocabulary,
named exemplars, evidence boundary, privacy class, protected objects, and
explicit unknowns.
\item[Argument blueprint] The claim map, evidence anchors, concept-dependency
graph, section jobs, terminology order, and visual plan approved before
long-form drafting.
\item[Document] The stable identity of a work, its genre, intended reader, and
declared decision.
\item[Source snapshot] The exact source tree, build environment, compiler
output, and content hash used for a version.
\item[Protected object] An equation, number, label, citation, table, quotation,
code block, or qualification with a source span and normalized form.
\item[Finding] A located passage, an observable reason for attention, an
optional machine suggestion, and the author's decision.
\item[Revision] A source snapshot derived from a parent version, with the
author, time, changed spans, and explanations.
\item[Reader session] A rendered version, reader role, context supplied, task
questions, time observations, and decision.
\item[Preflight run] The build, integrity, argument, writing, evidence,
reconstruction, mutation, abstention, and repair results for a candidate
version.
\item[Release decision] The gate results, explicit exceptions, packet hash,
unresolved risks, and named owner of release.
\item[Evidence item] A source or software observation tied to a claim, a
fixture, or a product choice, with provenance and appraisal state.
\end{description}

The relationships are more important than the field names. A brief owns a
blueprint; a blueprint owns bounded units; a finding belongs to a source
snapshot and a preflight run; a revision answers one or more findings; a
protected-object comparison links parent and child snapshots; and a reader
session sees only one released version at a time. An evidence item may motivate
a finding rule or a reader question, but it cannot silently alter a document.

Let $D$ be a document, $s$ a source snapshot, and $r$ a revision. The basic
lineage is
\begin{equation}
  D \longleftarrow B \longrightarrow P \longrightarrow s_0
  \xrightarrow{\;r_1\;} s_1 \longrightarrow \text{preflight}
  \longrightarrow \text{release} \longrightarrow \text{reader session}.
  \label{eq:architecture-lineage}
\end{equation}
The arrows carry a person and a time. A generated suggestion is not an arrow
until it is accepted into a bounded unit and the source changes; a candidate is
not a release until its preflight record passes. This simple rule makes
it possible to answer a later question such as ``who changed the denominator
in Table 3, and why?'' without reading a prompt transcript.

\section{The source snapshot and the canonical build}
\label{sec:architecture-snapshot}

When an author opens a version, the system must be able to reproduce the page
the reader will see. It therefore stores the whole source directory, not just a
hash of the top file: entry point, included files, bibliography, image hashes,
compiler version, package list, relevant environment variables, and build log.
The rendered PDF belongs to the version because readers judge the page, not the
source alone.

The build wrapper performs the checks in a fixed order. It first validates that
the requested entry point names one reader-facing route and that the source is a
read-only snapshot. It then compiles with \texttt{-no-shell-escape}, inside a
network-denied scratch directory, twice around bibliography generation, checks
references and citations, and finally records the PDF hash and page count. A
clean build is useful evidence about typography and reproducibility. It says
nothing about whether the argument is persuasive, so the wrapper reports it as
a build result rather than as a reader verdict. The complete trust-boundary
contract appears in Appendix~\ref{app:implementation-contract}.

The snapshot also stores a detexed prose view for diagnostics. Detexing is a
lossy transformation. Equations, tables, captions, and source locations remain
in the canonical source and PDF; the prose view is allowed to answer only
questions about words, sentence boundaries, and paragraph structure. A finding
created from the prose view links back to its source span before it can enter an
author's queue.

\begin{table}[H]
\centering
\small
\begin{tabularx}{\textwidth}{@{}p{0.23\textwidth}L p{0.28\textwidth}@{}}
\toprule
Snapshot field & Example value & Why the field matters \\
\midrule
Entry point & \path{docs/survey/humanvoice_survey.tex} & Establishes which
document the reader is being asked to judge \\
Included source hash & Hash over route files & Prevents an unrecorded include
from changing the evidence \\
Build environment & pdfTeX version, packages, date epoch & Makes a PDF
replayable and distinguishes source from rendering noise \\
Protected manifest & Object IDs, spans, normalized values & Makes a changed
number or equation visible \\
Authoring brief & Role, decision, time budget, disclosure, evidence boundary &
Defines the task and the material the writer may use \\
Blueprint hash & Claims, dependencies, unit jobs, visual plan & Makes the
argument order and writer context reproducible \\
Preflight and release & Gate results, mutations, abstentions, packet hash &
Shows why the reader was asked to spend time \\
\bottomrule
\end{tabularx}
\caption{A snapshot joins the source, its rendering, and the reader task.}
\label{tab:snapshot-fields}
\end{table}

\section{Protected objects and correspondence}
\label{sec:architecture-protection}

Before an author accepts a revision, the system must find every equation,
number, label, citation, table, quotation, code block, and qualification whose
silent alteration could change the argument. The extractor scans the source
and rendered representations, assigns each object a stable local ID, and keeps
the raw span beside a normalized value. It may ignore whitespace around an
equation operator; it may not ignore a changed exponent, sign, unit, or table
denominator.

For object $p$ in version $v$, let $raw_v(p)$ be the source span, $norm_v(p)$
the normalized representation, and $loc_v(p)$ its source and PDF location. A
candidate correspondence from parent $v$ to child $v'$ is accepted only if
the object type, semantic anchor, and local context agree:
\begin{equation}
  match(p_v,p_{v'}) =
  1\{type_v=type_{v'}\}\,1\{anchor_v\approx anchor_{v'}\}\,
  1\{context_v\approx context_{v'}\}.
  \label{eq:architecture-match}
\end{equation}
The relation $\approx$ is intentionally typed. For a citation it can use the
key and bibliography identity; for a table cell it can use row and column
labels; for an equation it can use its environment, label, and neighboring
text. When the relation is ambiguous, the extractor returns an unresolved
match for the author rather than choosing the highest textual similarity.

The comparison result has four ordinary-language outcomes: unchanged,
explained change, added or removed with an author explanation, and unresolved.
The last outcome blocks a release of the revision packet. It is the one place
where abstention is more valuable than coverage.

\begin{cardexample}{algebraic equivalence is not string equality}
The parent writes $a/(b+c)$ and the child writes $(b+c)^{-1}a$. A character
diff reports a large change. A parser may establish algebraic equivalence under
the declared commutativity assumptions, but it must also show the author the
assumptions and the source locations. If the parent instead writes
$a/(b-c)$, the same normalizer must not silently identify the two forms. The
first release can conservatively mark both changes for review; it need not
solve symbolic mathematics to protect a document.
\end{cardexample}

Numbers receive an analogous treatment. The normalized record stores value,
unit, sign, precision, and a relation to a table, equation, or prose claim.
Changing ``3.0 hours'' to ``3 hours'' is a formatting change if the precision
policy permits it. Changing ``3.0'' to ``30'' is a substantive change even when
the surrounding paragraph is identical. The user interface should show both
the raw and normalized forms so that a reader can tell which rule was applied.

\section{Findings as questions, not verdicts}
\label{sec:architecture-findings}

When a passage deserves attention, the author needs more than a score. The
finding layer combines deterministic rules, corpus statistics, and optional
model suggestions to provide a location, an observable reason, a consequence
hypothesis, and a next action. The hypothesis does not declare the passage
wrong; it explains why the passage may deserve a human's time.

For finding $f$, write
\begin{equation}
  f=(loc, rule, observation, consequence, question, suggestion),
  \label{eq:architecture-finding}
\end{equation}
where $suggestion$ may be empty. A useful finding can therefore read:

\begin{quote}
The phrase ``frozen interface'' occurs before the comparison is defined. A
reader without the project history may not know which inputs are held fixed.
Should the paragraph name the baseline and the changed mechanism first?
\end{quote}

The finding does not read ``AI phrase detected'' or ``bad writing.'' Those
labels collapse an observation, a genre judgment, and a remedy into one
authority claim. They also invite authors to optimize the label rather than
repair the passage.

Findings are ranked by expected reader consequence and review cost. A practical
priority function for the pilot is
\begin{equation}
  priority(f)=\frac{p_f\,c_f}{1+b_f},
  \label{eq:architecture-priority}
\end{equation}
where $p_f$ is the adjudicated probability that the issue blocks the reader's
task, $c_f$ is its consequence if left unrepaired, and $b_f$ is the expected
author burden. The terms are estimates used to order a queue, not a quality
score. The pilot should compare this ordering with the order in which readers
actually report problems.

\section{Three views with different responsibilities}
\label{sec:architecture-views}

The author view is a source-aware workbench. It shows the passage, source
location, rendered context, protected objects nearby, and any model suggestion.
The author can accept, reject, defer, or rewrite. A rejection asks for a short
reason, not a long compliance form. The view should make the next useful edit
obvious without forcing the author to inspect every low-confidence rule hit.

The reviewer view is narrower. An independent reviewer sees the source claim,
the cited evidence, the change explanation, and the protected comparison. The
reviewer can confirm the correspondence, ask for a correction, or mark the
claim as not checked. The reviewer does not need the author's entire prompt
history.

The reader view is deliberately plain. It is the rendered proposal, its
figures and references, and a short decision context. It does not show finding
IDs, model confidence, parser warnings, or internal phase names. At the end of
the reading task, the reader records the decision, the first unresolved
question, and the time spent. This separation is a product requirement: a
reader who sees the backstage machinery is no longer independent of the project.

\begin{table}[H]
\centering
\small
\begin{tabularx}{\textwidth}{@{}p{0.19\textwidth}p{0.28\textwidth}L@{}}
\toprule
View & Primary question & Information deliberately omitted \\
\midrule
Author & What should I investigate or change, and what meaning must I protect? &
Other readers' private judgments and irrelevant rule hits \\
Reviewer & Is the source, change, and evidence correspondence defensible? &
Prompt transcripts and a global style score \\
Reader & Can I understand the decision and act on this document? & Source
history, finding IDs, model confidence, and internal project language \\
\bottomrule
\end{tabularx}
\caption{The same revision has three views because the three people have
different jobs.}
\label{tab:architecture-views}
\end{table}

\section{Local inference and the trust boundary}
\label{sec:architecture-trust}

The product is local-first because technical documents often contain private
data, unpublished results, or contractual material. Local-first does not mean
that every computation must run on a laptop. It means that the default path
keeps the source inside an explicitly declared boundary and requires an
authorization before a remote service sees content.

The inference adapter records model family, model file or endpoint, prompt
template, sampling parameters, network state, and output hash. It passes only
the context needed for the finding. A citation identity lookup may use a public
identifier without sending the surrounding manuscript. A model suggestion may
be generated on a redacted passage. The author decides whether the suggestion
is useful; the adapter never writes directly to the canonical source.

Privacy is a property of the data flow. Let $X$ be the source, $Z$ a redacted
context, and $Y$ a model response. A local run has $X\to Z\to Y$ inside the
declared boundary. A remote run adds an external recipient $E$ and requires
explicit authorization:
\begin{equation}
  X\to Z\to Y\quad\text{(local)},\qquad
  X\to Z\to E\to Y\quad\text{(remote, authorized only)}.
  \label{eq:architecture-trust-flow}
\end{equation}
The diagram is a data-flow statement, not a promise that redaction is perfect.
The privacy review must test whether equations, rare identifiers, and quoted
sentences can re-identify a source even after obvious names are removed.

\section{Interfaces and commands}
\label{sec:architecture-interfaces}

The first command surface has the same six verbs introduced in
Chapter~\ref{ch:product} and specified in Chapter~\ref{ch:proposal}:

\begin{description}[style=nextline,leftmargin=30mm,labelindent=0mm]
\item[\texttt{hv init}] Validates the authoring brief, evidence boundary,
privacy class, and protected-object declarations. Critical blanks block the
writer.
\item[\texttt{hv plan}] Builds the claim map, dependency graph, unit jobs,
terminology order, and visual plan. Missing warrants become questions.
\item[\texttt{hv draft}] Invokes the replaceable writer for one approved unit
and records the context and evidence it received.
\item[\texttt{hv preflight}] Runs build, integrity, argument, writing,
evidence, reconstruction, and mutation checks. Unknown results remain
blockers.
\item[\texttt{hv repair}] Applies a named repair, rebuilds, reruns preflight,
and stops on a cycle limit or oscillation.
\item[\texttt{hv release}] Presents a clean version and records the named
reader's task outcome only after the release gate passes. A report is an export
of this record, not another top-level command.
\end{description}

The commands exchange versioned JSON records, not hidden global state. A
minimal finding record can be represented as
\begin{verbatim}
{
  "finding_id": "F-014",
  "source": "chapter.tex:87",
  "rule": "private-register",
  "observation": "phase label precedes public definition",
  "question": "name the comparison before the project history",
  "author_decision": "accepted",
  "revision_id": "R-003"
}
\end{verbatim}
The record schemas and compatibility rule are fixed in Appendix~\ref{app:implementation-contract}
and in \path{schemas/}. A minor schema version may add optional fields; a major
change requires migration and fixture replay. The invariant is that a later
reader can locate the passage and the author decision without opening an
orchestration log.

\section{Failure handling and observability}
\label{sec:architecture-failure}

The system should fail in ways that help a person decide what to do next. An
unknown macro yields ``source not parsed'' with a location. A missing citation
identifier yields ``identity unresolved'' with the key. A protected equation
that cannot be matched yields ``comparison unresolved'' and blocks the release
gate. A local model that exceeds the memory envelope yields a clear fallback to
deterministic findings, but any gate that depends on the model remains
abstained. A failed dependency or reconstruction check yields a located repair,
not a reader packet. None of these states should be phrased as a successful
review with a small confidence score.

When a parser cannot read a macro, the build log should say so. When a model
falls outside the memory envelope, the inference adapter should record the
fallback. The build log records compiler errors and warnings, the parser
records unsupported constructs, the adapter records latency, memory, and
network policy, and the reader protocol records task completion and questions.
Those records are what we mean by observability. A single ``quality'' number
would hide the failed boundary and make repair harder.

The operational envelope for the feasibility build is intentionally modest and
is stated numerically in Appendix~\ref{app:implementation-contract}: a 300-page
cap, 2,000 protected objects, bounded scratch space, and fixed deterministic
and model-assisted time limits on the reference machine. Those are planning
caps, not performance claims. Week six records the actual distributions and
either revises the caps openly or narrows the supported document class.

\section{The result of this architecture}
\label{sec:architecture-conclusion}

Follow one document through the system and the promise becomes testable. The
brief identifies the reader and decision; the blueprint explains why each unit
exists; a bounded draft records its evidence; preflight exposes a failure before
release; a revision records a human choice; a protected comparison exposes a
meaning-bearing change; and a clean reader view supplies the outcome. Engineers
can replace a component as long as it respects those boundaries.

The architecture does not guarantee a persuasive document. It makes the
failure legible and makes a comparative experiment possible. That is the
appropriate ambition for the first investment. The next Part specifies the
experiment, the economic decision rule, and the operating work required if the
path earns a larger trial.
